% \documentclass[tikz, border=10pt]{standalone}
% \usepackage{amsmath}
% \usetikzlibrary{positioning, arrows.meta, fit, backgrounds, calc}

% \begin{document}
% \begin{tikzpicture}[
%     node distance = 7mm and 12mm, % Increased horizontal spacing to prevent overlaps
%     font=\small,
%     every node/.style = {align=center},
%     box/.style = {
%         rectangle, rounded corners, draw=black, thick,
%         minimum width=3.3cm, minimum height=1.1cm,
%         text width=3.1cm, align=center
%     },
%     plain/.style = {
%         rectangle, rounded corners, draw=black!60,
%         minimum width=3.3cm, minimum height=0.9cm,
%         text width=3.1cm, align=center, fill=white
%     },
%     blueBox/.style    = {box, fill=blue!20},
%     orangeBox/.style = {box, fill=orange!25},
%     greenBox/.style  = {box, fill=green!20},
%     grayBox/.style   = {box, fill=gray!20},
%     branchTitle/.style = {font=\small\bfseries, align=center, text width=3.6cm},
%     solidArrow/.style = {-{Latex[length=2.2mm]}, thick},
%     dashArrow/.style  = {-{Latex[length=2.2mm]}, thick, dashed, draw=black!70},
% ]

% %% ---------------- Stage 1: Input (Far Left) ----------------
% \node[plain] (qc) {Quantum Circuit};
% \node[plain, right=of qc] (cfe) {Circuit Feature\\Encoding};
% \node[blueBox, right=of cfe] (fe) {Feature\\Embedding};

% %% ---------------- Horizontal Row Coordinates ----------------
% % Fixed Y positions relative to the Feature Embedding node to isolate the branches
% \coordinate (upperRowY) at ($(fe.center) + (0,  2.2cm)$);
% \coordinate (lowerRowY) at ($(fe.center) + (0, -2.5cm)$);

% %% ---------------- Stage 2A: Original NNAS branch (Upper Row) ----------------
% \node[branchTitle] (titleL) at ($(fe.east) + (2.2cm, 2.2cm)$) {Original NNAS\\(Stochastic Branch)};
% \node[blueBox, right=of titleL] (na)   {Neural\\Accumulator};
% \node[blueBox, right=of na]         (sls) {Stochastic\\Latent State};

% %% ---------------- Stage 2B: Proposed coherent branch (Lower Row) ----------------
% \node[branchTitle] (titleR) at ($(fe.east) + (2.2cm, -2.5cm)$) {Proposed Physics-Informed\\Coherent Branch};
% \node[orangeBox, right=5mm of titleR] (re) {Recurrent\\Encoder};
% \node[orangeBox, right=of re]         (pn) {Posterior Network\\Outputs: $\mu_l,\ \log\sigma_l^2$};
% \node[orangeBox, right=of pn]         (rp) {Reparameterization};
% \node[orangeBox, right=of rp]         (sd) {Sample $\delta h_l$};

% %% ---------------- Gate-conditioned prior (Placed safely above posterior) ----------------
% \node[grayBox, above=9mm of pn] (gp) {Gate-conditioned\\Prior Network\\$p(\delta h_l \mid g_l)$};

% %% ---------------- Stage 3: Physics-based coherent propagation (Green Boxes Expanded) ----------------
% \node[greenBox, right=of sd] (bch) {Physics-Informed\\BCH Propagation\\$\Delta H_{l-1} + \mathrm{Ad}_U(\delta h_l)$};
% \node[greenBox, right=of bch] (dhl) {Accumulated\\Coherent Generator\\$\Delta H_l$};
% \node[font=\scriptsize\itshape, below=1mm of bch] (bchnote) {Deterministic};

% %% ---------------- Stage 4: Fusion & Convergence ----------------
% \node[blueBox, right=of dhl] (proj) {Projection\\of $\Delta H_l$};

% % Dynamic center anchor for fusion between upper and lower rows
% \node[blueBox] (fus) at ($(proj.east |- fe.center) + (1.5cm, 0)$) {Latent Fusion\\(Concatenation)};

% %% ---------------- Stage 5: Prediction ----------------
% \node[blueBox, right=of fus] (ne) {Neural\\Extractor};
% \node[blueBox, right=of ne] (mh) {Mitigation\\Head};
% \node[blueBox, right=of mh] (po) {Predicted\\Observable};

% %% ---------------- Training losses (Far Right Layout Fix) ----------------
% \node[grayBox, above right=2mm and 12mm of po] (ml)  {Mitigation Loss};
% \node[grayBox, below right=2mm and 12mm of po] (kll) {KL Loss};
% \node[grayBox, right=4.5cm of po] (tl)  {Total Loss};

% %% ================= ARROWS: main inference path =================
% \draw[solidArrow] (qc) -- (cfe);
% \draw[solidArrow] (cfe) -- (fe);

% \draw[solidArrow] (fe.east) -- ++(4mm,0) |- (titleL.west);
% \draw[solidArrow] (fe.east) -- ++(4mm,0) |- (titleR.west);

% \draw[solidArrow] (titleL) -- (na);
% \draw[solidArrow] (na) -- (sls);

% \draw[solidArrow] (titleR) -- (re);
% \draw[solidArrow] (re) -- (pn);
% \draw[solidArrow] (pn) -- (rp);
% \draw[solidArrow] (rp) -- (sd);
% \draw[solidArrow] (sd) -- (bch);
% \draw[solidArrow] (bch) -- (dhl);
% \draw[solidArrow] (dhl) -- (proj);

% \draw[solidArrow] (sls.east) -- ++(8mm,0) |- (fus.west);
% \draw[solidArrow] (proj.east) -- ++(8mm,0) |- (fus.west);

% \draw[solidArrow] (fus) -- (ne);
% \draw[solidArrow] (ne) -- (mh);
% \draw[solidArrow] (mh) -- (po);

% \draw[solidArrow] (po.east) -- ++(4mm,0) |- (ml.west);
% \draw[solidArrow] (po.east) -- ++(4mm,0) |- (kll.west);
% \draw[solidArrow] (ml.east) -- ++(5mm,0) |- (tl.west);
% \draw[solidArrow] (kll.east) -- ++(5mm,0) |- (tl.west);

% %% ================= ARROWS: training-only (dashed) =================
% \draw[dashArrow] (gp.south) -- node[midway, right, font=\scriptsize] {KL Reg.} (pn.north);

% % Clear pathing around the right side of the graphics grid for KL calculation
% \draw[dashArrow] (gp.north) -- ++(0, 5mm) -| (kll.north);
% \draw[dashArrow] (pn.south) -- ++(0,-7mm) -| ($(re.south east) + (8mm, -5mm)$) |- (kll.south);

% %% ================= Legend (Bottom Left placement) =================
% \begin{scope}[shift={($(qc.south) + (-0.5cm, -4.5cm)$)}]
%     \node[blueBox, minimum width=1.0cm, text width=0cm, minimum height=0.4cm] (legB) {};
%     \node[right=1mm of legB, font=\scriptsize, align=left] {Original NNAS modules};
    
%     \node[orangeBox, minimum width=1.0cm, text width=0cm, minimum height=0.4cm, right=4.0cm of legB] (legO) {};
%     \node[right=1mm of legO, font=\scriptsize, align=left] {New generative latent modules};
    
%     \node[greenBox, minimum width=1.0cm, text width=0cm, minimum height=0.4cm, right=4.5cm of legO] (legG) {};
%     \node[right=1mm of legG, font=\scriptsize, align=left] {Physics-informed deterministic modules};
    
%     \node[grayBox, minimum width=1.0cm, text width=0cm, minimum height=0.4cm, right=5.3cm of legG] (legY) {};
%     \node[right=1mm of legY, font=\scriptsize, align=left] {Training-only components};
% \end{scope}

% \end{tikzpicture}
% \end{document}

\documentclass[tikz, border=10pt]{standalone}
\usepackage{amsmath}
\usetikzlibrary{positioning, arrows.meta, fit, backgrounds, calc}

\begin{document}
\begin{tikzpicture}[
    node distance = 7mm and 10mm,
    font=\small,
    every node/.style = {align=center},
    box/.style = {
        rectangle, rounded corners, draw=black, thick,
        minimum width=3.2cm, minimum height=1.05cm,
        text width=3.0cm, align=center
    },
    plain/.style = {
        rectangle, rounded corners, draw=black!60,
        minimum width=3.2cm, minimum height=0.9cm,
        text width=3.0cm, align=center, fill=white
    },
    blueBox/.style   = {box, fill=blue!20},
    orangeBox/.style = {box, fill=orange!25},
    greenBox/.style  = {box, fill=green!20},
    grayBox/.style   = {box, fill=gray!20},
    branchTitle/.style = {font=\small\bfseries, align=center, text width=3.6cm},
    solidArrow/.style = {-{Latex[length=2.2mm]}, thick},
    dashArrow/.style  = {-{Latex[length=2.2mm]}, thick, dashed, draw=black!70},
]

%% ---------------- Stage 1: Input ----------------
\node[plain] (qc) {Quantum Circuit};
\node[plain, below=of qc] (cfe) {Circuit Feature\\Encoding};
\node[blueBox, below=of cfe] (fe) {Feature\\Embedding};

%% ---------------- Branch anchors ----------------
% Left column (original NNAS) and right column (proposed coherent branch),
% symmetric about the Feature Embedding node.
\coordinate (splitL) at ($(fe.south) + (-3.0cm, -1.0cm)$);
\coordinate (splitR) at ($(fe.south) + ( 3.0cm, -1.0cm)$);

\node[branchTitle, below=3mm of splitL] (titleL) {Original NNAS\\(Stochastic Branch)};
\node[branchTitle, below=3mm of splitR] (titleR) {Proposed Physics-Informed\\Coherent Branch};

%% ---------------- Stage 2A: Original NNAS branch (left) ----------------
\node[blueBox, below=8mm of titleL] (na)  {Neural\\Accumulator};
\node[blueBox, below=of na]         (sls) {Stochastic\\Latent State};

%% ---------------- Stage 2B: Proposed coherent branch (right) ----------------
\node[orangeBox, below=8mm of titleR] (re) {Recurrent\\Encoder};
\node[orangeBox, below=of re]         (pn) {Posterior Network\\Outputs: $\mu_l,\ \log\sigma_l^2$};
\node[orangeBox, below=of pn]         (rp) {Reparameterization};
\node[orangeBox, below=of rp]         (sd) {Sample $\delta h_l$};

%% ---------------- Gate-conditioned prior (gray, training-only) ----------------
\node[grayBox, right=14mm of pn] (gp) {Gate-conditioned\\Prior Network\\$p(\delta h_l \mid g_l)$};

%% ---------------- Stage 3: Physics-based coherent propagation (green) ----------------
\node[greenBox, below=of sd] (bch) {Physics-Informed\\BCH Propagation\\$\Delta H_{l-1} + \mathrm{Ad}_U(\delta h_l)$};
\node[font=\scriptsize\itshape, below=0.5mm of bch] (bchnote) {Deterministic};
\node[greenBox, below=9mm of bchnote] (dhl) {Accumulated\\Coherent Generator\\$\Delta H_l$};

%% ---------------- Stage 4: Fusion ----------------
\node[blueBox, below=of dhl] (proj) {Projection\\of $\Delta H_l$};

\coordinate (fusY) at ($(proj.south) + (0, -1.1cm)$);
\node[blueBox] (fus) at ($(fusY)!0.5!(sls |- fusY)$) {Latent Fusion\\(Concatenation)};

%% ---------------- Stage 5: Prediction ----------------
\node[blueBox, below=of fus] (ne) {Neural\\Extractor};
\node[blueBox, below=of ne] (mh) {Mitigation\\Head};
\node[blueBox, below=of mh] (po) {Predicted\\Observable};

%% ---------------- Training losses ----------------
\node[grayBox, below=13mm of po, xshift=-2.1cm] (ml)  {Mitigation Loss};
\node[grayBox, below=13mm of po, xshift= 2.1cm] (kll) {KL Loss};
\node[grayBox, below=of po, yshift=-2.6cm] (tl)  {Total Loss};

%% ================= ARROWS: main inference path =================
\draw[solidArrow] (qc) -- (cfe);
\draw[solidArrow] (cfe) -- (fe);

\draw[solidArrow] (fe.south) -- ++(0,-3mm) -| (na.north);
\draw[solidArrow] (fe.south) -- ++(0,-3mm) -| (re.north);

\draw[solidArrow] (na) -- (sls);
\draw[solidArrow] (re) -- (pn);
\draw[solidArrow] (pn) -- (rp);
\draw[solidArrow] (rp) -- (sd);
\draw[solidArrow] (sd) -- (bch);
\draw[solidArrow] (bch) -- (dhl);
\draw[solidArrow] (dhl) -- (proj);

\draw[solidArrow] (sls.south) |- ($(fus.west)+(0,0)$) -- (fus.west);
\draw[solidArrow] (proj.south) -- ++(0,-3mm) -| (fus.north);

\draw[solidArrow] (fus) -- (ne);
\draw[solidArrow] (ne) -- (mh);
\draw[solidArrow] (mh) -- (po);

\draw[solidArrow] (po.south) -- ++(0,-3mm) -| (ml.north);
\draw[solidArrow] (ml) -- (tl);
\draw[solidArrow] (kll) -- (tl);

%% ================= ARROWS: training-only (dashed) =================
% Prior regularizes the posterior (KL regularization), drawn near the top.
\draw[dashArrow] (gp.north) to[out=140, in=0]
    node[midway, above, font=\scriptsize, text width=2.2cm, align=center] {KL\\Regularization}
    (pn.east);

% Prior and posterior both feed the KL loss at the bottom; routed down the
% right margin of the figure (clean orthogonal paths) to avoid crossing the
% main inference path.
\coordinate (marginX) at ($(gp.east) + (2.0cm, 0)$);
\draw[dashArrow] (gp.east) -| (marginX) |- (kll.east);
\draw[dashArrow] (pn.east) -| ($(marginX) + (0.6cm,0)$) |- (kll.north);

%% ================= Legend =================
% Anchored to the top-right of the figure (relative to Feature Embedding)
% rather than to a mid-figure node, so it never collides with the branch
% content below.
\begin{scope}[shift={($(fe.east)+(6.5cm,1.0cm)$)}]
    \node[blueBox, minimum width=1.6cm, text width=0cm, minimum height=0.5cm] (legB) {};
    \node[right=1mm of legB, font=\scriptsize, text width=4cm, align=left] {Original NNAS modules};
    \node[orangeBox, minimum width=1.6cm, text width=0cm, minimum height=0.5cm, below=2mm of legB] (legO) {};
    \node[right=1mm of legO, font=\scriptsize, text width=4cm, align=left] {New generative latent modules};
    \node[greenBox, minimum width=1.6cm, text width=0cm, minimum height=0.5cm, below=2mm of legO] (legG) {};
    \node[right=1mm of legG, font=\scriptsize, text width=4cm, align=left] {Physics-informed deterministic modules};
    \node[grayBox, minimum width=1.6cm, text width=0cm, minimum height=0.5cm, below=2mm of legG] (legY) {};
    \node[right=1mm of legY, font=\scriptsize, text width=4cm, align=left] {Training-only components};
\end{scope}

\end{tikzpicture}
\end{document}